Fix \(p\in\mathfrak O\) and define its latent fiber
\[
F(p)=\{v\in\Delta_{\mathrm{SLCC}}:Bv=p\}.
\]
It is nonempty because \(p\in B\Delta_{\mathrm{SLCC}}\), convex because its constraints are linear, and compact because it is a closed subset of the finite-dimensional compact simplex.  By \(\text{thm:finite-response-converse}\), the compatible causal values are exactly
\[
\{\tau_e(M):M\in\mathfrak M_{\mathrm{SLCC}},\ P_M(O_i)=p\}
=\{h_e^{\top}v:v\in F(p)\}.
\tag{1}
\]
The continuous linear image on the right of (1) is a nonempty compact convex subset of \(\mathbb R\), hence a compact interval.  Its attained minimum and maximum are \(L(p)\) and \(U(p)\) by \(\text{def:sharp-endpoint-programs}\).

To verify that no value between the endpoints is omitted, choose minimizer \(v_L\) and maximizer \(v_U\).  For \(0\leq\gamma\leq1\),
\[
v_\gamma=(1-\gamma)v_L+\gamma v_U\in F(p),
\qquad
h_e^{\top}v_\gamma=(1-\gamma)L(p)+\gamma U(p).
\]
The converse construction in \(\text{thm:finite-response-converse}\) produces a legal SCM for every \(v_\gamma\).  Thus the causal identified set is exactly \([L(p),U(p)]\), both endpoints are attained, and every intermediate value is attained.